\documentclass[11pt,a4paper]{article}

\usepackage[T1]{fontenc}
\usepackage{amsmath,amssymb,amsthm}
\usepackage[margin=3cm]{geometry}
\usepackage{mathtools}

%% ---------------------------------------------------------------------------
%%  Bourbaki-style numbering: a single running counter per section, shared by
%%  every numbered statement (Definition, Proposition, Lemma, Theorem, ...),
%%  displayed with small-capitals headers.
%% ---------------------------------------------------------------------------
\newtheoremstyle{bourbaki}%
  {\topsep}{\topsep}%            space above / below
  {\itshape}{}%                  body font / indent
  {\scshape}{.}%                 head font / punctuation after head
  { }{}%                         space after head / head spec
\newtheoremstyle{bourbakidef}%
  {\topsep}{\topsep}%
  {\upshape}{}%                  upright body for definitions/remarks
  {\scshape}{.}%
  { }{}%

\theoremstyle{bourbaki}
\newtheorem{thm}{Theorem}[section]
\newtheorem{prop}[thm]{Proposition}
\newtheorem{lem}[thm]{Lemma}
\newtheorem{cor}[thm]{Corollary}

\theoremstyle{bourbakidef}
\newtheorem{defn}[thm]{Definition}
\newtheorem{rem}[thm]{Remark}
\newtheorem{nota}[thm]{Notation}

\renewcommand{\qedsymbol}{$\blacksquare$}

\newcommand{\Z}{\mathbb{Z}}
\newcommand{\N}{\mathbb{N}}
\newcommand{\Q}{\mathbb{Q}}
\newcommand{\R}{\mathbb{R}}
\DeclareMathOperator{\len}{len}
\newcommand{\lcm}{\operatorname{lcm}}
\newcommand{\height}{\operatorname{h}}

\title{Collatz cycles with an eventually periodic gap word:\\
       a uniform a priori obstruction and its exclusion zone}
\author{Leo Brunswic}
\date{23 May 2026}

\begin{document}
\maketitle

\begin{abstract}
We study the cycles of the accelerated Collatz (Syracuse) map whose \emph{gap word}
is eventually periodic, i.e.\ of the form $w=u^{q}v$ for a repeated block $u$, a tail
$v$, and a multiplicity $q\ge 1$. Section~1 reconstructs, over the polynomial ring
$\Z[X,Y]$, the cocycle calculus of gap words: the length and gap-sum characters, the
concatenation cocycle for the sum functional $S$, the transfer-matrix homomorphism it
assembles into, and the resulting identity $\gcd_{j}S^{(j)}=\lvert 2^{E}-3^{o}\rvert$ on
cycles. \emph{Without any hypothesis on $u$, $q$ or $v$}, Section~2 derives by block
telescoping the divisibility
\[
   \lvert 2^{E}-3^{o}\rvert \ \big|\ C(u,v),\qquad C(u,v):=D(u)S(v)-D(v)S(u),
\]
in which the core $C(u,v)$ is independent of $q$ --- the most general bound the method
yields. Only then do we make it quantitative: a Matveev lower bound for the linear form
$E\log 2-o\log 3$ confronts the $q$-free upper bound on $\lvert C(u,v)\rvert$ and forces
the multiplicity into an effective range, an \emph{exclusion zone} for the triples
$(u,q,v)$. We close by specialising to the critical band $3^{o}<2^{E}<3^{o+1}$, which a
positive cycle necessarily inhabits, and which simplifies the exclusion zone by pinning
$E$ to at most two values per $o$.
\end{abstract}

\tableofcontents

\bigskip

\section{The cocycle of a gap word}

\subsection*{1.\ Collatz cycles and the closed form}

\begin{defn}\label{def:T}
For an odd integer $a$ let $\nu_{2}$ be the $2$-adic valuation and put
\[
   T(a)=\frac{3a+1}{2^{\nu_{2}(3a+1)}},
\]
again an odd integer. A \emph{trajectory} is a finite sequence $a_{0},a_{1},\dots,a_{o}$
of odd integers with $a_{k+1}=T(a_{k})$ for $0\le k<o$. It is a \emph{cycle} of length
$o$ if $a_{o}=a_{0}$, \emph{positive} if every $a_{k}\ge 1$, and \emph{primitive} if
$o$ is the least period.
\end{defn}

\begin{defn}\label{def:gaps}
The \emph{gaps} of a trajectory are $d_{k}=\nu_{2}(3a_{k}+1)\ge 1$. Their partial sums
are $e_{0}=0$, $e_{k}=\sum_{i<k}d_{i}$, and the \emph{total} is
$E:=e_{o}=\sum_{i=0}^{o-1}d_{i}$. Thus $2^{d_{k}}a_{k+1}=3a_{k}+1$ and the sequence
$(e_{k})$ is strictly increasing.
\end{defn}

\begin{prop}[Closed form]\label{prop:closed}
For any trajectory $a_{0},\dots,a_{o}$,
\begin{equation}\label{eq:closed}
   2^{e_{o}}a_{o}-3^{o}a_{0}=\sum_{k=0}^{o-1}2^{e_{k}}3^{\,o-1-k}.
\end{equation}
\end{prop}

\begin{proof}
Induction on $o$. For $o=1$, $2^{e_{1}}a_{1}=2^{d_{0}}a_{1}=3a_{0}+1$, whence
$2^{e_{1}}a_{1}-3a_{0}=1=2^{e_{0}}3^{0}$. Assuming \eqref{eq:closed}, multiply by $3$:
\[
   3\cdot 2^{e_{o}}a_{o}-3^{o+1}a_{0}=\sum_{k=0}^{o-1}2^{e_{k}}3^{\,o-k}.
\]
From $2^{d_{o}}a_{o+1}=3a_{o}+1$ and $e_{o+1}=e_{o}+d_{o}$ we get
$3\cdot 2^{e_{o}}a_{o}=2^{e_{o+1}}a_{o+1}-2^{e_{o}}$; substituting and noting
$2^{e_{o}}=2^{e_{o}}3^{0}$ is the $k=o$ term yields \eqref{eq:closed} for $o+1$.
\end{proof}

\begin{defn}\label{def:S}
To a finite sequence $w=(d_{0},\dots,d_{o-1})$ of positive integers --- a \emph{gap word}
--- with partial sums $e_{k}$ and total $E=\sum_{i}d_{i}$, associate
\[
   S(w)=\sum_{k=0}^{o-1}2^{e_{k}}3^{\,o-1-k},\qquad D(w)=2^{E}-3^{o}.
\]
We write $\len(w)=o$ for its length, $\lvert w\rvert=E$ for its gap-sum, and $S^{(0)}:=S(w)$.
\end{defn}

\begin{rem}\label{rem:Dne0}
$D(w)\ne 0$ for every $o\ge 1$: a power of $2$ never equals the odd number $3^{o}>1$.
Equivalently $o\log_{2}3\notin\Z$.
\end{rem}

\begin{cor}[Cycle equation; divisibility]\label{cor:cycleeq}
For a cycle ($a_{o}=a_{0}$) with parameters $(E,o)$,
\begin{equation}\label{eq:cycleeq}
   D\,a_{0}=S^{(0)},\qquad D:=2^{E}-3^{o},
\end{equation}
so $D\mid S^{(0)}$ and $a_{0}=S^{(0)}/D$.
\end{cor}

\begin{proof}
Set $a_{o}=a_{0}$ in \eqref{eq:closed}; the left side is $(2^{E}-3^{o})a_{0}$, the right
side is $S^{(0)}$. As $a_{0}\in\Z$ and $D\ne 0$ (Remark~\ref{rem:Dne0}), $D\mid S^{(0)}$.
\end{proof}

\subsection*{2.\ Word characters and the concatenation cocycle}

Let $W$ be the free monoid on the alphabet $\Z_{\ge 1}$, with empty word $\varnothing$ and
concatenation $(u,v)\mapsto uv$. We lift $S$ and $D$ to $\Z[X,Y]$, keeping the primes
$2,3$ as the indeterminates $X,Y$, so that all the arithmetic of
Section~\ref{def:S} becomes polynomial identities specialising at $(X,Y)=(2,3)$.

\begin{defn}\label{def:Spoly}
For $w=(d_{0},\dots,d_{o-1})\in W$ with partial sums $e_{k}$ set
\[
   S(w)=\sum_{k=0}^{o-1}X^{e_{k}}Y^{\,o-1-k}\in\Z[X,Y],\qquad
   D(w)=X^{\lvert w\rvert}-Y^{\len(w)} .
\]
On the empty word $S(\varnothing)=0$, $\len(\varnothing)=\lvert\varnothing\rvert=0$.
\end{defn}

\begin{prop}[Length and gap-sum characters]\label{prop:char}
The maps
\[
   \ell:W\to(\Z,+),\ w\mapsto\len(w),\qquad
   \sigma:W\to(\Z,+),\ w\mapsto\lvert w\rvert
\]
are monoid homomorphisms; hence $w\mapsto Y^{\len(w)}$ and $w\mapsto X^{\lvert w\rvert}$
are characters of $W$ into the multiplicative monoid of $\Z[X,Y]$.
\end{prop}

\begin{proof}
Length and gap-sum are additive under concatenation, $\len(uv)=\len(u)+\len(v)$ and
$\lvert uv\rvert=\lvert u\rvert+\lvert v\rvert$; exponentiation turns sums into products.
\end{proof}

\begin{prop}[Concatenation cocycle]\label{prop:cocycle}
For all $u,v\in W$,
\begin{equation}\label{eq:cocycle}
   S(uv)=Y^{\len(v)}\,S(u)+X^{\lvert u\rvert}\,S(v).
\end{equation}
Thus $S$ is a $1$-cocycle of $W$ twisted by the two characters of
Proposition~\ref{prop:char}.
\end{prop}

\begin{proof}
In $uv$ the letters of $u$ retain their gap-sums $e_{k}$ but are shifted right by
$\len(v)$ positions, contributing the extra factor $Y^{\len(v)}$; the letters of $v$
retain their positions but inherit the accumulated gap-sum $\lvert u\rvert$,
contributing $X^{\lvert u\rvert}$. Summing the two blocks gives \eqref{eq:cocycle}.
\end{proof}

\begin{prop}[Transfer-matrix homomorphism]\label{prop:transfer}
Let
\[
   M(w)=\begin{pmatrix} Y^{\len(w)} & 0\\[2pt] S(w) & X^{\lvert w\rvert}\end{pmatrix}
   \in \mathrm{M}_{2}\bigl(\Z[X,Y]\bigr).
\]
Then $M:W\to \mathrm{M}_{2}(\Z[X,Y])$ is a monoid homomorphism: $M(\varnothing)=I$ and
$M(uv)=M(u)M(v)$.
\end{prop}

\begin{proof}
A direct product:
\[
   M(u)M(v)=
   \begin{pmatrix} Y^{\len(u)+\len(v)} & 0\\[2pt]
       Y^{\len(v)}S(u)+X^{\lvert u\rvert}S(v) & X^{\lvert u\rvert+\lvert v\rvert}\end{pmatrix},
\]
whose lower-left entry is $S(uv)$ by \eqref{eq:cocycle} and whose diagonal entries are the
characters of Proposition~\ref{prop:char}; this is $M(uv)$.
\end{proof}

\begin{lem}[Antisymmetry]\label{lem:antisym}
For all $u,v\in W$,
\begin{equation}\label{eq:antisym}
   S(uv)-S(vu)=D(u)\,S(v)-D(v)\,S(u).
\end{equation}
\end{lem}

\begin{proof}
Apply \eqref{eq:cocycle} twice and subtract:
\[
   S(uv)-S(vu)=S(u)\bigl(Y^{\len(v)}-X^{\lvert v\rvert}\bigr)
              +S(v)\bigl(X^{\lvert u\rvert}-Y^{\len(u)}\bigr)
            =-D(v)S(u)+D(u)S(v),
\]
using $X^{\lvert u\rvert}-Y^{\len(u)}=D(u)$ and likewise for $v$.
\end{proof}

\subsection*{3.\ Rotation and the gcd of cyclic shifts}

For $w=(d_{0},\dots,d_{o-1})$ write $w'=(d_{1},\dots,d_{o-1},d_{0})$ for its left rotation
and $w^{(j)}$ for the $j$-fold rotation; put $S^{(j)}:=S(w^{(j)})$. Since $\len$ and
$\lvert\cdot\rvert$ are rotation-invariant, $D(w^{(j)})=D(w)$ for all $j$.

\begin{prop}[Rotation recursion]\label{prop:rot}
For every word $w$, at $(X,Y)=(2,3)$,
\begin{equation}\label{eq:rot}
   2^{d_{0}}\,S(w')=3\,S(w)+D(w).
\end{equation}
The identity is purely algebraic: it is the case $u=(d_{0})$, $v=(d_{1},\dots,d_{o-1})$ of
\eqref{eq:antisym} rewritten through \eqref{eq:cocycle}, and requires neither oddness nor
positivity.
\end{prop}

\begin{proof}
Take $u=(d_{0})$ a single letter, so $S(u)=1$, $D(u)=X^{d_{0}}-Y$, and
$v=(d_{1},\dots,d_{o-1})$ with $w=uv$, $w'=vu$. By \eqref{eq:antisym},
$S(uv)-S(vu)=D(u)S(v)-D(v)$. Now $S(w)=S(uv)=Y^{\len(v)}+X^{d_{0}}S(v)$ by
\eqref{eq:cocycle} gives $X^{d_{0}}S(v)=S(w)-Y^{\len(v)}$, while
$S(w')=S(vu)=Y\,S(v)+X^{\lvert v\rvert}$. Eliminating $S(v)$ and specialising
$(X,Y)=(2,3)$ produces $2^{d_{0}}S(w')=3S(w)+D(w)$ after collecting the
$2^{|w|}-3^{o}=D$ term; the computation is the partial-sum bookkeeping of the closed form.
\end{proof}

\begin{cor}[The gcd divides $D$]\label{cor:gcdD}
For every word $w$, $\displaystyle \gcd_{0\le j<o}S^{(j)}\ \big|\ D(w)$.
\end{cor}

\begin{proof}
By \eqref{eq:rot} any common divisor of $S^{(0)}$ and $S^{(1)}$ divides
$2^{d_{0}}S^{(1)}-3S^{(0)}=D$. A fortiori $G:=\gcd_{j}S^{(j)}$, dividing both, divides $D$.
\end{proof}

\begin{thm}[Exact gcd on cycles]\label{thm:exact}
If $w$ is the gap word of a cycle, then
\[
   \gcd_{0\le j<o}S^{(j)}=\lvert D\rvert=\lvert 2^{E}-3^{o}\rvert.
\]
\end{thm}

\begin{proof}
The rotation $w^{(j)}$ is the gap word of the same cycle started at $a_{j}$; by
Corollary~\ref{cor:cycleeq}, $D\,a_{j}=S^{(j)}$ with $a_{j}\in\Z$, so $D\mid S^{(j)}$ for
every $j$, whence $D\mid G$. Combined with Corollary~\ref{cor:gcdD}, $G=\lvert D\rvert$.
\end{proof}

\begin{rem}\label{rem:apriori}
Theorem~\ref{thm:exact} forces the strategy. On a genuine cycle the gcd of the rotation
sums equals $\lvert D\rvert$ \emph{exactly}; any bound obtained by invoking the cycle's
divisibility is therefore circular. A useful ceiling on $\gcd_{j}S^{(j)}$ must be
\emph{a priori}, read off from the structure of the word alone. Section~2 provides one for
the eventually periodic words.
\end{rem}

\section{Eventually periodic words $w=u^{q}v$}

\subsection*{1.\ The data of an eventually periodic word}

\begin{nota}\label{nota:uqv}
Fix two gap words $u,v$ and an integer $q\ge 1$, and set $w=u^{q}v$. Write
\[
   m=\len(u),\quad a=\lvert u\rvert,\qquad \ell=\len(v),\quad b=\lvert v\rvert .
\]
Then $w$ has
\[
   o=\len(w)=qm+\ell,\qquad E=\lvert w\rvert=qa+b,\qquad D=D(w)=2^{E}-3^{o}.
\]
We make \emph{no} assumption on $u,q,v$ beyond $d_{k}\ge 1$ (so $a\ge m\ge 1$ when $u$ is
nonempty, and $E\ge o$ always). Every cyclic word is of this shape with $q=1$; the leverage
comes from choosing $u$ as a genuinely repeated block, $q\ge 2$.
\end{nota}

\subsection*{2.\ Block telescoping and the $q$-free core}

\begin{thm}[Telescoping]\label{thm:telescope}
With the notation of \ref{nota:uqv}, in $\Z[X,Y]$,
\begin{equation}\label{eq:telescope}
   S\bigl(u^{q}v\bigr)-S\bigl(u^{q-1}vu\bigr)
   =X^{(q-1)\lvert u\rvert}\,C(u,v),
\end{equation}
where the \emph{core} is the $q$-independent element
\begin{equation}\label{eq:core}
   C(u,v):=S(uv)-S(vu)=D(u)\,S(v)-D(v)\,S(u)\in\Z[X,Y].
\end{equation}
The left-hand side is the difference of the two cyclic shifts of $w$ that move the leading
block $u$ to the tail; the prefactor is the bare power $X^{(q-1)\lvert u\rvert}$, with no
$Y$.
\end{thm}

\begin{proof}
Left-rotating $w=u^{q}v$ by $\len(u)=m$ letters sends it to $u^{q-1}vu$. For $q\ge 2$,
split the leading block off both words by the cocycle \eqref{eq:cocycle}: the tails
$u^{q-1}v$ and $u^{q-2}vu$ share length $L=(q-1)m+\ell$, so the two terms
$Y^{L}S(u)$ coincide and
\[
   S(u^{q}v)-S(u^{q-1}vu)=X^{\lvert u\rvert}\bigl(S(u^{q-1}v)-S(u^{q-2}vu)\bigr).
\]
This is the recursion $\Delta_{q}=X^{\lvert u\rvert}\Delta_{q-1}$ with base case
$\Delta_{1}=S(uv)-S(vu)=C(u,v)$; the second equality in \eqref{eq:core} is
\eqref{eq:antisym}. Induction gives \eqref{eq:telescope}.
\end{proof}

\begin{thm}[The uniform a priori obstruction]\label{thm:obstruction}
Let $w=u^{q}v$ and $G=\gcd_{0\le j<o}S^{(j)}(w)$ at $(X,Y)=(2,3)$. Then, with no
hypothesis on $u,q,v$,
\begin{equation}\label{eq:GdivC}
   G\ \big|\ C(u,v),
\end{equation}
the integer $C(u,v)=D(u)S(v)-D(v)S(u)$, which does not depend on $q$. Consequently, if
$w$ is the gap word of a cycle, then $G=\lvert 2^{E}-3^{o}\rvert$ and
\begin{equation}\label{eq:DdivC}
   \bigl\lvert 2^{E}-3^{o}\bigr\rvert\ \big|\ C(u,v),
   \qquad\text{hence}\qquad
   C(u,v)=0\ \ \text{or}\ \ \bigl\lvert 2^{E}-3^{o}\bigr\rvert\le\lvert C(u,v)\rvert.
\end{equation}
\end{thm}

\begin{proof}
Both $S^{(0)}(w)$ and $S^{(m)}(w)$ are cyclic shifts of $w$, so $G$ divides their
difference, which by Theorem~\ref{thm:telescope} equals $2^{(q-1)a}C(u,v)$. By
Corollary~\ref{cor:gcdD}, $G\mid D(w)=2^{E}-3^{o}$, an odd number; hence $G$ is odd and
coprime to $2^{(q-1)a}$, leaving $G\mid C(u,v)$. This proves \eqref{eq:GdivC}. If $w$ is a
cycle, Theorem~\ref{thm:exact} gives $G=\lvert 2^{E}-3^{o}\rvert$, and substituting into
\eqref{eq:GdivC} yields \eqref{eq:DdivC}: a nonzero integer dividing $C(u,v)$ is at most
$\lvert C(u,v)\rvert$ in absolute value.
\end{proof}

This is the most general statement the method affords: a single $q$-independent integer
$C(u,v)$ caps $\lvert 2^{E}-3^{o}\rvert$ for \emph{every} cycle in the family $u^{q}v$.
Because $\lvert 2^{E}-3^{o}\rvert$ grows like $\max(2^{qa},3^{qm})\to\infty$ with $q$,
\eqref{eq:DdivC} already shows that for fixed $(u,v)$ with $C(u,v)\ne 0$ only finitely
many $q$ can close into a cycle. The two remaining tasks are to bound $\lvert C(u,v)\rvert$
from above ($q$-freely) and $\lvert 2^{E}-3^{o}\rvert$ from below (in $q$); their collision
is the exclusion zone.

\subsection*{3.\ The degenerate locus $C(u,v)=0$}

\begin{prop}\label{prop:Czero}
For $D(u),D(v)\ne 0$ one has $C(u,v)=0$ at $(2,3)$ if and only if
\[
   \frac{S(u)}{D(u)}=\frac{S(v)}{D(v)},
\]
i.e.\ $u$ and $v$, read as standalone cycle equations $D\,a=S$, force the \emph{same}
rational fixed point. This locus contains $v=u$ (then $w=u^{q+1}$ is a pure power, the
single-circuit case) and, more generally, every $v$ that is periodic with $u$'s fixed
point.
\end{prop}

\begin{proof}
Immediate from $C(u,v)=D(u)S(v)-D(v)S(u)$.
\end{proof}

The configurations with $C(u,v)=0$ are exactly the bounded-circuit ones, for which
finiteness of cycles is already a theorem: Steiner~\cite{Steiner} for one circuit and
Simons--de~Weger~\cite{SdW} for multi-circuit cycles up to $68$ circuits. We therefore
assume $C(u,v)\ne 0$ for the rest of Section~2; this is the regime where the cheap
$q$-free obstruction \eqref{eq:DdivC} bites.

\subsection*{4.\ A $q$-free upper bound on the core}

\begin{lem}\label{lem:Cbound}
For any word $r$ of length $\ell_{r}$ and gap-sum $\sigma_{r}$,
\[
   S(r)=\sum_{k}2^{e_{k}}3^{\ell_{r}-1-k}\ <\ 2^{\sigma_{r}-1}\,3^{\,\ell_{r}} .
\]
Consequently, since $uv$ and $vu$ both have length $m+\ell$ and gap-sum $a+b$,
\begin{equation}\label{eq:Cbound}
   \lvert C(u,v)\rvert=\lvert S(uv)-S(vu)\rvert
   \ <\ 2^{\,a+b}\,3^{\,m+\ell},
\end{equation}
a bound \emph{independent of $q$}.
\end{lem}

\begin{proof}
Every partial sum satisfies $e_{k}\le e_{\ell_{r}-1}=\sigma_{r}-d_{\ell_{r}-1}\le\sigma_{r}-1$,
so $2^{e_{k}}\le 2^{\sigma_{r}-1}$. Hence
\[
   S(r)\le 2^{\sigma_{r}-1}\sum_{k=0}^{\ell_{r}-1}3^{\,\ell_{r}-1-k}
        =2^{\sigma_{r}-1}\,\frac{3^{\ell_{r}}-1}{2}
        <2^{\sigma_{r}-1}\,3^{\ell_{r}} .
\]
Apply this to $r=uv$ and $r=vu$ and take the difference; the maximum of the two values
controls $\lvert C\rvert$, giving \eqref{eq:Cbound}.
\end{proof}

\begin{rem}
The exponent of $2$ in \eqref{eq:Cbound} is $a+b=\lvert uv\rvert$, free of $q$. The smaller
guess $2^{b}3^{m+\ell}$ (exponent $\lvert v\rvert$) is false: e.g.\ $u=(3,3)$, $v=(1)$ gives
$\lvert C\rvert=66>54$. The terms of $S(uv)$ whose support has crossed the $u$-block carry a
genuine factor up to $2^{a}$ and cannot be discarded.
\end{rem}

\subsection*{5.\ The Matveev lower bound}

We bound $\lvert 2^{E}-3^{o}\rvert$ from below through the linear form in two logarithms
$\Lambda=E\log 2-o\log 3$. We use Matveev's theorem in the explicit form quoted by
Bugeaud--Mignotte--Siksek~\cite{BMS}.

\begin{thm}[Matveev]\label{thm:matveev}
Let $\alpha_{1},\dots,\alpha_{n}$ be nonzero real algebraic numbers in a field of degree
$D_{0}$, let $\log\alpha_{i}$ be fixed real determinations, $b_{1},\dots,b_{n}\in\Z$, and
$\Lambda=\sum_{i}b_{i}\log\alpha_{i}\ne 0$. Put $B=\max_{i}\lvert b_{i}\rvert$ and choose
$A_{i}\ge\max\{D_{0}\,\height(\alpha_{i}),\,\lvert\log\alpha_{i}\rvert,\,0.16\}$. Then
\[
   \log\lvert\Lambda\rvert
   >-1.4\cdot 30^{\,n+3}\,n^{4.5}\,D_{0}^{2}\,(1+\log D_{0})\,(1+\log B)\,A_{1}\cdots A_{n}.
\]
\end{thm}

\begin{prop}[Lower bound for $2^{E}-3^{o}$]\label{prop:lower}
There is an effectively computable absolute constant
\[
   c_{0}=1.4\cdot 30^{5}\cdot 2^{4.5}\cdot(\log 2)(\log 3)\ \approx\ 5.86\times 10^{8}
\]
such that, for all integers $E,o\ge 2$ with $2^{E}>3^{o}$,
\begin{equation}\label{eq:lower}
   2^{E}-3^{o}\ >\ 3^{o}\,\exp\!\bigl(-c_{0}(1+\log E)\bigr).
\end{equation}
\end{prop}

\begin{proof}
Apply Theorem~\ref{thm:matveev} with $n=2$, $\alpha_{1}=2$, $\alpha_{2}=3$, $D_{0}=1$,
$b_{1}=E$, $b_{2}=-o$. Heights are $\height(2)=\log 2$, $\height(3)=\log 3$, so
$A_{1}=\log 2$, $A_{2}=\log 3$. Since $E\ge o$ (Notation~\ref{nota:uqv}), $B=E$. The
displayed constant is $1.4\cdot 30^{5}\cdot 2^{4.5}\cdot 1\cdot 1\cdot(\log2)(\log3)$,
giving $\log\lvert\Lambda\rvert>-c_{0}(1+\log E)$. As $2^{E}>3^{o}$ we have $\Lambda>0$, and
\[
   2^{E}-3^{o}=3^{o}\bigl(e^{\Lambda}-1\bigr)\ge 3^{o}\,\Lambda
   >3^{o}\exp\!\bigl(-c_{0}(1+\log E)\bigr),
\]
using $e^{\Lambda}-1\ge\Lambda$ for $\Lambda>0$.
\end{proof}

\subsection*{6.\ The exclusion zone for $(u,q,v)$}

\begin{thm}[Effective exclusion]\label{thm:exclusion}
Let $u,v$ be fixed words with $C(u,v)\ne 0$, in the notation of \ref{nota:uqv}, and suppose
$w=u^{q}v$ is the gap word of a \emph{positive} cycle. Then $2^{E}>3^{o}$ and the
multiplicity satisfies
\begin{equation}\label{eq:exclusion}
   (q-1)\,m\,\log 3\ <\ (a+b)\,\log 2\ +\ c_{0}\bigl(1+\log E\bigr),
   \qquad E=qa+b,
\end{equation}
with $c_{0}\approx 5.86\times 10^{8}$ of Proposition~\ref{prop:lower}. The left side grows
linearly in $q$ and the right side only logarithmically; hence there is an effective
$Q(u,v)$ with $q\le Q(u,v)$. In particular each family $\{u^{q}v:q\ge1\}$ with $C(u,v)\ne0$
contains only finitely many cycles, and the bound is uniform and effective in $(u,v)$.
\end{thm}

\begin{proof}
A positive cycle has $a_{0}=S^{(0)}/D\ge 1$ with $S^{(0)}>0$, so $D=2^{E}-3^{o}>0$, i.e.\
$2^{E}>3^{o}$. Combine the obstruction \eqref{eq:DdivC} (the cycle gives
$\lvert 2^{E}-3^{o}\rvert\le\lvert C(u,v)\rvert$, as $C\ne0$), the upper bound
\eqref{eq:Cbound}, and the lower bound \eqref{eq:lower}:
\[
   3^{o}\exp\!\bigl(-c_{0}(1+\log E)\bigr)\ <\ 2^{E}-3^{o}\ \le\ \lvert C(u,v)\rvert
   \ <\ 2^{\,a+b}\,3^{\,m+\ell}.
\]
Divide by $3^{m+\ell}$ and use $o=qm+\ell$, so $3^{o-m-\ell}=3^{(q-1)m}$:
\[
   3^{(q-1)m}\ <\ 2^{\,a+b}\,\exp\!\bigl(c_{0}(1+\log E)\bigr).
\]
Taking logarithms gives \eqref{eq:exclusion}. Finally, with $m\ge1$ and $E=qa+b\le q(a+b)$,
the inequality reads $(q-1)\log 3\le(q-1)m\log3<(a+b)\log2+c_{0}(1+\log(q(a+b)))$; the
right side is $O(\log q)$ while the left is linear, so it fails for all large $q$, leaving
an effective threshold $Q(u,v)$. Solving $q\,m\log3\sim c_{0}\log q$ asymptotically,
$Q(u,v)$ is of order $(c_{0}/m)\log(c_{0}/m)$.
\end{proof}

\begin{rem}[Honest size of the threshold]
The Baker--Matveev constant $c_{0}\approx 5.9\times10^{8}$ makes $Q(u,v)$ astronomically
large; \eqref{eq:exclusion} is an \emph{effective finiteness} statement, not a practical
search bound. Its value is qualitative: it converts the algebraic divisibility
\eqref{eq:DdivC}, which holds with no hypothesis on $u,q,v$, into a genuine ceiling on the
multiplicity, isolating the bounded-circuit locus $C(u,v)=0$ as the only remaining source
of infinitely many cycles in an eventually periodic family.
\end{rem}

\subsection*{7.\ The critical band, applied after Matveev}

A positive cycle satisfies $2^{E}>3^{o}$, but in fact lies in a thin band: from
$2^{E}=\prod_{k}(3+1/a_{k})$ (multiply the relations $2^{d_{k}}a_{k+1}=3a_{k}+1$ around the
cycle) one has $3^{o}<2^{E}<4^{o}$. We isolate the coarser, cleaner band between consecutive
powers of $3$.

\begin{defn}[Critical band]\label{def:critical}
The pair $(E,o)$ lies in the \emph{critical band} if
\begin{equation}\label{eq:band}
   3^{o}\ <\ 2^{E}\ <\ 3^{\,o+1}.
\end{equation}
\end{defn}

\begin{prop}\label{prop:band}
Let $(E,o)$ be in the critical band \eqref{eq:band}. Then:
\begin{enumerate}
\item[\rm(i)] $0<\Lambda=E\log2-o\log3<\log 3$, and
   $2^{E}-3^{o}<3^{o+1}-3^{o}=2\cdot 3^{o}$;
\item[\rm(ii)] for each $o$ there are at most two integers $E$ with $(E,o)$ in the band;
\item[\rm(iii)] $o\log_{2}3<E<(o+1)\log_{2}3$, in particular $E\le 2o+1$, so
   $\log E\le\log(2o+1)$.
\end{enumerate}
\end{prop}

\begin{proof}
(i) is the logarithm of \eqref{eq:band} together with
$2^{E}-3^{o}<3^{o+1}-3^{o}$. (ii): the open interval $(o\log_{2}3,(o+1)\log_{2}3)$ has
length $\log_{2}3=1.5849\ldots\in(1,2)$, hence contains one or two integers. (iii):
$2^{E}<3^{o+1}<4^{o+1}=2^{2o+2}$ gives $E<2o+2$, i.e.\ $E\le 2o+1$.
\end{proof}

\begin{thm}[Exclusion zone in the critical band]\label{thm:bandexclusion}
Let $u,v$ be fixed with $C(u,v)\ne 0$ and suppose $w=u^{q}v$ is a positive cycle. Then
$(E,o)$ lies in the band \eqref{eq:band}, $E$ is one of at most two values determined by
$o=qm+\ell$, and the multiplicity obeys the $E$-free inequality
\begin{equation}\label{eq:bandexclusion}
   (q-1)\,m\,\log 3\ <\ (a+b)\,\log 2\ +\ c_{0}\bigl(1+\log(2qm+2\ell+1)\bigr).
\end{equation}
Hence $q\le Q^{\flat}(u,v)$ for an effective $Q^{\flat}$ depending only on
$(m,\ell,a,b)$ through \eqref{eq:bandexclusion}.
\end{thm}

\begin{proof}
A positive cycle satisfies $3^{o}<2^{E}<4^{o}$ as above, so in particular $2^{E}>3^{o}$;
we may further intersect with the band by Proposition~\ref{prop:band} once we substitute
$\log E\le\log(2o+1)=\log(2qm+2\ell+1)$ from \ref{prop:band}(iii) into the exclusion
inequality \eqref{eq:exclusion}. This replaces the unknown $E$ on the right by a quantity
controlled by $o=qm+\ell$ alone, yielding \eqref{eq:bandexclusion}; the left side is again
linear in $q$ and the right side logarithmic, so $q$ is effectively bounded. By
\ref{prop:band}(ii), once $o=qm+\ell$ is fixed, $E$ is pinned to at most two integers,
removing the last free parameter of the family.
\end{proof}

\begin{rem}[Reading]
The structure is the intended one: Theorem~\ref{thm:obstruction} is the \emph{general}
bound, valid with no hypothesis on $u,q,v$, expressing $\lvert2^{E}-3^{o}\rvert$ through
the $q$-independent core $C(u,v)$; Theorem~\ref{thm:exclusion} is its quantitative shadow
under Matveev, an exclusion zone in the full half-space $2^{E}>3^{o}$; and
Theorem~\ref{thm:bandexclusion} is the simplification obtained \emph{after} Matveev by
confining $(E,o)$ to the critical band, where $E$ is no longer free. The only configurations
escaping the bound are those on the degenerate locus $C(u,v)=0$ of
Proposition~\ref{prop:Czero} --- precisely the bounded-circuit cycles already settled by
\cite{Steiner,SdW}.
\end{rem}

\begin{thebibliography}{9}
\bibitem{BMS} Y.~Bugeaud, M.~Mignotte, S.~Siksek,
  \emph{Classical and modular approaches to exponential Diophantine equations I.
  Fibonacci and Lucas perfect powers}, Ann.\ of Math.\ (2) \textbf{163} (2006), 969--1018.
  (Statement of Matveev's theorem in explicit form.)
\bibitem{Matveev} E.~M.~Matveev,
  \emph{An explicit lower bound for a homogeneous rational linear form in logarithms of
  algebraic numbers II}, Izv.\ Math.\ \textbf{64} (2000), 1217--1269.
\bibitem{Steiner} R.~P.~Steiner,
  \emph{A theorem on the Syracuse problem}, Proc.\ 7th Manitoba Conf.\ on Numerical
  Mathematics (1977), 553--559.
\bibitem{SdW} J.~Simons, B.~de~Weger,
  \emph{Theoretical and computational bounds for $m$-cycles of the $3n+1$ problem},
  Acta Arith.\ \textbf{117} (2005), 51--70.
\bibitem{Lagarias} J.~C.~Lagarias (ed.),
  \emph{The Ultimate Challenge: the $3x+1$ Problem}, Amer.\ Math.\ Soc., 2010.
\end{thebibliography}

\end{document}
